\section{数据采购与自建：买数据还是造数据}

产业里常见选择是买数据、采数据、造数据。买数据速度快但同质化，自建数据壁垒深但成本高，合成/仿真数据可扩展但有真实性风险。不同阶段选择不同：早期验证可以购买或合成，进入核心能力后必须建立自有闭环。

\subsection{三种路径对比}

\noindent\begin{tabularx}{\textwidth}{p{0.20\textwidth}p{0.34\textwidth}X}
\toprule
路径 & 适用阶段 & 风险 \\
\midrule
买数据 & 快速启动、补充通用能力 & 同质化、授权不清、边际收益低。 \\
采数据 & 建立真实任务能力 & 成本高、周期长、隐私和安全复杂。 \\
造数据 & 长尾、危险、稀缺场景 & 分布偏移、评测失真、过拟合仿真。 \\
\bottomrule
\end{tabularx}

\begin{importantbox}{长期壁垒来自自有闭环}
可以买来的数据很难成为长期护城河。真正的壁垒来自自有任务、自有反馈、自有评测和自有 Recipe 的组合。
\end{importantbox}

\subsection{本章小结}

数据策略不是买或造的二选一。短期可以买和合成，长期必须形成自己的数据闭环和评测体系。

